\chapter{Operational distinctions}\label{chap:C}

The following distinctions guide the connected case analysis. They were used to describe what changed, who acted and how the activity continued. They do not constitute exhaustive event-frequency coding. [\hyperlink{src:D44}{D44}, \hyperlink{src:D21}{D21}]

The examples in Table~\ref{tab:C-2} show how the distinctions were applied to connected sequences. The actor, piece, local change and source time are carried forward separately from an explanation of motive. Participant numbers are stable across conditions; the original software endpoints are linked through the checked recording mappings.

\begingroup
\setstretch{1}
\fontsize{9.5}{11.685}\selectfont
\setlength{\tabcolsep}{2.2pt}
\renewcommand{\arraystretch}{1.08}
\begin{longtable}{@{}>{\raggedright\arraybackslash}p{28.439mm}>{\raggedright\arraybackslash}p{62.763mm}>{\raggedright\arraybackslash}p{65.705mm}@{}}
\caption[Operational distinctions used for the case analysis.]{Operational distinctions used for the case analysis.}\label{tab:C-1}\\
\toprule
\textbf{Distinction} & \textbf{Evidence used} & \textbf{Interpretive boundary} \\
\midrule
\endfirsthead
\multicolumn{3}{l}{\small\itshape Table \thetable{} (continued)}\\
\toprule
\textbf{Distinction} & \textbf{Evidence used} & \textbf{Interpretive boundary} \\
\midrule
\endhead
\midrule
\multicolumn{3}{r}{\small Continued on next page}\\
\endfoot
\bottomrule
\endlastfoot
Reference work & A piece, location, orientation or target relation is identified or clarified. & A common camera orientation does not establish the same interpretation. \\
\addlinespace[3pt]
Local uptake & A response or placement follows a specific proposal or existing arrangement. & The sequence concerns that local step, not necessarily a shared whole solution. \\
\addlinespace[3pt]
Parallel and complementary work & Partners develop different parts, continue one another's placements or combine compatible contributions. & The allocation can change within a round; it is not a fixed assigned role. \\
\addlinespace[3pt]
Self-revision & The person who proposed or built a local arrangement later changes or dismantles it. & Verbal self-correction and physical reconfiguration are recorded separately. \\
\addlinespace[3pt]
Partner revision & A partner changes an existing placement or introduces an actionable alternative. & The action alone does not establish a motive for intervening. \\
\addlinespace[3pt]
Preservation and tidying & An existing local arrangement is left in place or adjusted while its structure is retained. & Preservation by itself does not establish agreement, politeness or lack of alternatives. \\
\addlinespace[3pt]
Assistance and manipulation difficulty & An observed or reported difficulty is followed by a request, adjustment or partner action. & The sequence does not isolate viewpoint from control, practice or task demands. \\
\addlinespace[3pt]
Intervention and outcome & Researcher feedback, declared completion, on-site acceptance and later geometric checking are distinguished. & A local coordinated action and a correct final silhouette are separate judgments. \\
\end{longtable}
\endgroup

\begingroup
\setstretch{1}
\fontsize{9.5}{11.685}\selectfont
\setlength{\tabcolsep}{2.2pt}
\renewcommand{\arraystretch}{1.18}
\begin{longtable}{@{}>{\raggedright\arraybackslash}p{29.420mm}>{\raggedright\arraybackslash}p{63.744mm}>{\raggedright\arraybackslash}p{63.744mm}@{}}
\caption[Worked examples of the operational distinctions, retaining source-specific timestamps.]{Worked examples of the operational distinctions, retaining source-specific timestamps.}\label{tab:C-2}\\
\toprule
\textbf{Episode and source} & \textbf{Observed or reported sequence} & \textbf{Applied distinction and supported interpretation} \\
\midrule
\endfirsthead
\multicolumn{3}{l}{\small\itshape Table \thetable{} (continued)}\\
\toprule
\textbf{Episode and source} & \textbf{Observed or reported sequence} & \textbf{Applied distinction and supported interpretation} \\
\midrule
\endhead
\midrule
\multicolumn{3}{r}{\small Continued on next page}\\
\endfoot
\bottomrule
\endlastfoot
G01 R02, goose head and neck. \hyperlink{src:D03}{D03}; researcher report at 04:38, 06:01 and 06:26. & P01 joined the small triangle and parallelogram as a head/neck at 04:38. P02 tidied the two-piece arrangement at 06:01. P01 rotated the parallelogram towards a large triangle at 06:26. The researcher judged the first local arrangement correct and the revised one incorrect. & Preservation/tidying followed by self-revision. The original builder later changed the local solution. The account identifies actors, materials and sequence; it does not specify why P02 retained it or prove that tidying consisted only of edge alignment. \\
\addlinespace[3pt]
G03 R03, alternative flame constructions. \hyperlink{src:D06}{D06}, P05 video 03:06; \hyperlink{src:D21}{D21}, active task seconds. & The researcher reported P05 dismantling the two-small-triangle flame and reusing its pieces for the base. Separately, the log records P05 holding S1 at 67.376--68.875 s and S3 at 69.329--71.514 s, followed by P06 holding S3 at 93.555--94.394 s and S1 at 96.291--97.708 s. & Self-revision and redistribution of material, followed by partner uptake. The video report establishes the local task meaning; the log establishes repeated use by two endpoints. The participant mapping was checked against visible grabs; video 03:06 and the selected log intervals retain their own time bases. \\
\addlinespace[3pt]
G06 R04, requested rotation assistance. \hyperlink{src:D15}{D15}, \hyperlink{src:D14}{D14} about 20:24; \hyperlink{src:D21}{D21}, software R03. & The researcher reported P11 requesting that P12 rotate and place the parallelogram. In the retained overhead trace, S4 changes from three P11 holds at 146.932--154.210 s to three P12 holds at 155.237--161.167 s. The earlier aborted attempt remains separate. & Requested assistance and a change of operator for the same piece. The source account provides the request; the log supports the endpoint sequence. The approximate video anchor does not establish a millisecond response time, and the handover does not by itself identify who originated the solution. \\
\end{longtable}
\endgroup

Video mm:ss and log active seconds are separate time bases. Values rounded here are navigation and sequence locators, not response-time measurements. The examples illustrate use of the distinctions, not an exhaustive frequency coding.

